\clearpage
\thispagestyle{empty}
\pdfbookmark[0]{Edition and sources}{translation-sources}
\begin{center}\textbf{Edition and Sources}\end{center}
\noindent This English translation contains the historical title pages, the preface, the contents, and all 192 articles of Carl Friedrich Gauss’s \textit{Theoria motus corporum coelestium in sectionibus conicis solem ambientium} (Hamburg, 1809), as printed in his \textit{Werke}, vol.~VII (Göttingen and Leipzig, 1906), together with the three tables, four figures, and Brendel’s historical remarks. It also contains all 18 Additions to the Theoria motus and the historical remarks on them. It also includes the nine notes and letters on cometary orbits and Gauss’s New Table of Parabolic Motion with its complete explanation and all three parts. The letters to Encke, Olbers, and the editor, and Brendel’s remarks conclude the parabolic-motion section. The section on the perturbations of Ceres includes both methods, the letters and tables, and Brendel’s remarks. The section on Pallas contains all 37 letters from the 1906 edition and the complete \textit{Exposition of a New Method for Calculating Planetary Perturbations}, including both drafts of its fourth section. The calculations include both accounts of the special perturbations by Jupiter, the first general calculation, and the second general calculation through the first fragment numbered [11] on p. 561 of the 1906 edition. The second fragment numbered [11], on pp. 561–564, is outside this edition’s scope. The tables for Jupiter’s perturbations and the sections on the perturbations by Saturn and Mars follow. Brendel’s remarks accompany the Pallas material. The lunar theory comprises all five sections and Brendel’s remarks. The volume closes with the historical corrections and additions, F. Klein’s note, the historical contents and printer’s imprint, the unnumbered blank leaf, and the publisher’s catalogue. The prices of 25 Mk. (A. IV), 30 Mk. (A. VII), and 180 Mk. (B, Vols. I–VI) follow the qualified readings of the 1906 copy. 
\par\medskip\noindent The historical text follows the copy digitized by the \href{https://gdz.sub.uni-goettingen.de/id/PPN236008730}{Göttinger Digitalisierungszentrum}. The unsigned title vignette uses the same design as the 1906 title leaf. Its reproduction is from the \href{https://commons.wikimedia.org/w/index.php?oldid=1179960952}{1870 second printing of Volume I}, University of California Libraries copy, marked \href{https://creativecommons.org/publicdomain/mark/1.0/}{Public Domain Mark 1.0}.
\par\medskip\noindent AI translation from Latin, German, and French. AI typesetting.
